\section{Normal induction and exact-field source replacement}
\label{sec:SF-normal-induction-descent}

\paragraph{Status and provenance.}
Accepted after complete author, root and independent whole-proof
readings; exact source scopes and documentary bindings are recorded.
The question of minimality and the warning about kernels on different
inducing blocks were raised by /root.  The Galois-orbit replacement and
the resulting normal-induction bounds were derived by /root/sixfolds;
the equal-kernel repair was independently identified by both agents.
The original-field and polarized-Frobenius inputs below are the accepted
D-OG-01--03 and HD-PD01, not new claims of their independent origin.

\subsection*{The actual source and the statement}

Let $K$ be imaginary quadratic.  Let $C\to\mathbf P^1$ be an actual
connected finite $G$-Galois cover in a family, and let $\rho_K$ be a
faithful absolutely irreducible $K$-representation of $G$ whose exact
complex character field is $K$.  Suppose its multiplicity variation has
$K$-dimension six, signature $(3,3)$, and actual period rank nine.
Its naturally polarized primitive factor is obtained by a
Rosati-fixed rank-one idempotent in the corresponding rational group
algebra block.  Fix its signed determinant class
$\delta\in\mathbf Q^\times/N_{K/\mathbf Q}(K^\times)$.
The theorem preserves this class; it does not assert equality of
integral polarization types.

Suppose
\[
 H\triangleleft G,\qquad
 \rho_{\mathbf C}\simeq\operatorname{Ind}_H^G\tau,
 \qquad e=\dim_{\mathbf C}\tau\ge2,
 \qquad k=[G:H],
 \tag{NID.1}
\]
where $\tau$ is absolutely irreducible.  Primitivity of $\tau$ is not
needed.  Put $N=\ker\tau$, $J_\tau=H/N$, $L=\mathbf Q(\chi_\tau)$,
and $\ell=[L:K]$.

\begin{theorem}[D-TD56: normal-induction replacement]
There is a subgroup $H\subset H'\subset G$ and an absolutely
irreducible exact-$K$ model $W_K$ such that
\[
 [H':H]=\ell,\qquad H'/H\simeq\operatorname{Gal}(L/K),
 \qquad W_{\mathbf C}\simeq\operatorname{Ind}_H^{H'}\tau,
 \qquad \operatorname{Ind}_{H'}^G W_K\simeq\rho_K.
 \tag{NID.2}
\]
The kernel of $W_K$ is exactly $N$.  In particular $N$ is normal in
$H'$, although it need not be normal in the original $G$.
The actual connected cover
\[
 Y=C/N\longrightarrow Z=C/H'
 \quad\hbox{has group }J'=H'/N,
 \qquad |J'|=\ell|J_\tau|,\qquad \dim_K W_K=e\ell.
 \tag{NID.3}
\]
It supplies an exact-$K$ source of signature $(3,3)$, actual period
rank nine, and the same natural signed determinant $\delta$.
Both $Z$ and $X=C/H$ have genus zero.

If $\ell>1$, the abelian group $H'/H$ is cyclic or a Klein four
group.  In the cyclic case
\[
       \ell\le\left\lfloor\frac{e+3}{5}\right\rfloor.
 \tag{NID.4}
\]
In the Klein-four case $\ell=4$ and $e\ge15$.  Consequently, with
\[
 b(e)=\max\left(4,\left\lfloor\frac{e+3}{5}\right\rfloor\right),
 \tag{NID.5}
\]
every source in (NID.1) has a replacement satisfying
\[
       \dim_K W_K\le e\,b(e),\qquad |J'|\le b(e)|J_\tau|.
 \tag{NID.6}
\]
If the original coefficient degree is minimal among all actual
exact-$K$, signature-$(3,3)$, rank-nine sources of natural determinant
class $\delta$, including nonprimitive sources, then $k=\ell$.
In that case $H'=G$, $N=1$, and the same bounds apply to the original
coefficient degree and original finite group order.
\end{theorem}

\subsection*{1. The regular Clifford orbit and its Galois action}

Since $H$ is normal, Mackey decomposition and Frobenius reciprocity give
\[
 1=\langle\rho,\rho\rangle_G
  =\sum_{gH\in G/H}\langle\tau,\tau^g\rangle_H.
 \tag{NID.7}
\]
Thus all $k$ conjugates $\tau^g$ are distinct, each occurs once in
$\operatorname{Res}_H^G\rho$, and the inertia group of $\tau$ is
exactly $H$.  The quotient $Q=G/H$ therefore acts regularly on their
set $\mathcal T$.

Induced-character values are sums of values of $\chi_\tau$; hence
$K\subset L$.  Character values lie in a cyclotomic field, whose
Galois group over $\mathbf Q$ is abelian.  Every subfield, including
$L$, is consequently Galois over $\mathbf Q$.  In particular
$\Gamma=\operatorname{Gal}(L/K)$ is an abelian group of order $\ell$.
Coefficientwise Galois action commutes with conjugation by $G$.
Because $\rho$ has a $K$-model, it permutes the constituents of its
restriction to $H$.  The action of $\Gamma$ on the orbit of $\tau$
is free: an automorphism fixes $\chi_\tau$ precisely when it fixes
the field generated by all of its values.  Thus
\[
        \mathcal O=\{\tau^\sigma:\sigma\in\Gamma\}
        \subset\mathcal T,\qquad |\mathcal O|=\ell.
 \tag{NID.8}
\]

The centralizer of a regular action of $Q$ is its opposite regular
action.  Here this elementary fact can be seen by choosing $\tau$ as
basepoint: any commuting permutation is determined by its value at
$\tau$, since its value at $q\tau$ must be $q$ times that value.
The elements of $Q$ carrying $\tau$ into $\mathcal O$ form a subgroup
$A$ isomorphic to $\Gamma^{\mathrm{op}}$.  Indeed, for $\sigma\in
\Gamma$ let $a_\sigma\tau=\tau^\sigma$; regularity gives a unique
$a_\sigma$, and commutation gives the opposite multiplication law.
Since $\Gamma$ is abelian, $A\simeq\Gamma$.
It is also the stabilizer of $\mathcal O$, and it acts regularly on
$\mathcal O$.  Let $H'$ be its preimage in $G$.  This proves the
subgroup and quotient assertions in (NID.2), without a choice of a
Galois lift as an automorphism of the whole cover.

\subsection*{2. An actual $K$-model and its exact kernel}

For an irreducible character $\chi$ of $H$ let $p_\chi$ be its
central complex group-algebra idempotent.  The orbit sum
\[
 p_{\mathcal O}=\sum_{\sigma\in\Gamma}p_{\chi_\tau^\sigma}
   =\frac{e}{|H|}\sum_{h\in H}
       \operatorname{Tr}_{L/K}\bigl(\chi_\tau(h^{-1})\bigr)h
       \ \in K[H]
 \tag{NID.9}
\]
is an idempotent.  Apply it to the given $K$-model of $\rho$ and put
$W_K=p_{\mathcal O}\rho_K$.  This is an actual $H'$-stable
$K$-subspace, not a hypothetical descent of one complex constituent.
Its complexification is the direct sum of the $\ell$ distinct
$H$-types in $\mathcal O$.  Since $H'$ acts transitively on these
multiplicity-one types, that complex representation is irreducible:
an invariant subspace, restricted to the semisimple $H$-module, must
be a sum of a subset of the types, and a nonempty invariant subset
is the whole orbit.  Frobenius reciprocity and dimensions identify it
with $\operatorname{Ind}_H^{H'}\tau$.

The $G$-translates of $W_K$ are the disjoint orbit sums attached to
the translates of $\mathcal O$.  They give the original $K$-module
and establish the last isomorphism in (NID.2).  Alternatively,
induction in stages gives that isomorphism after complexification,
and the nonzero intertwiner equations over $K$ give it over $K$.
The character field $K_0$ of $W_K$ lies in $K$.  Since induction
expresses every value of $\chi_\rho$ as a sum of values of
$\chi_{W_K}$, one also has $K\subset K_0$.  Thus $K_0=K$.
In particular Schur index one and the exact field are both supplied
by this construction; neither is inferred from character-field
equality alone.

All the characters $\tau^\sigma$ have the same kernel $N$.
For a finite-order complex matrix the character equals its degree
if and only if the matrix is the identity: after unitary
conjugation, equality in the triangle inequality forces all
eigenvalues to be one.  This criterion is unchanged by a field
automorphism.  On $H$, the kernel of $W_K$ is therefore exactly $N$.
If $h'\notin H$, its nontrivial image in $A$ regularly permutes the
nonzero constituent spaces, so it cannot act as the identity on
$W_K$.  This proves
\[
                       \ker(W_K)=N\subset H.
 \tag{NID.10}
\]
It proves normality in $H'$ and the asserted faithful quotient
$J'=H'/N$.  It does not assert normality of $N$ in $G$ before
minimality has forced $H'=G$.

\subsection*{3. The actual curve, period, and polarization}

Marking the finite subgroups after a finite cover of the parameter
space gives the actual family of quotient curves in (NID.3).
Frobenius reciprocity and finite quotient pullback identify its
multiplicity variation with the original one:
\[
 \operatorname{Hom}_G(\rho_K,H^1(C,K))
  \simeq \operatorname{Hom}_{H'}(W_K,H^1(C,K))
  \simeq \operatorname{Hom}_{J'}(W_K,H^1(Y,K)).
 \tag{NID.11}
\]
These are isomorphisms of variations induced by fixed group-algebra
correspondences.  They preserve the $K$-Hodge filtrations, so the
dimension, signature, and period rank on the actual original-family
image are unchanged.

For completeness, the rational polarization step is the exact
same-field argument of D-OG-02--03 and HD-PD01.  The rational block
of $W_K$ is $\operatorname{Mat}_{e\ell}(K)$.  Choose a
Rosati-fixed rank-one idempotent $p$ in it and lift it to $\mathbf
Q[H']$.  If $e_N=|N|^{-1}\sum_{n\in N}n$, the element $e_Np$
acts on $H^1(C,\mathbf Q)$.  It is supported on the $\rho$ block
only and has $K$-rank one there: every other irreducible containing
$W$ on restriction would occur in $\operatorname{Ind}_{H'}^G W$,
which is exactly $\rho$, and its occurrence has multiplicity one.
The quotient pullback $Y\leftarrow C$ identifies the selected
primitive factors and multiplies their polarizations by $|N|$.

Within the $\rho$ block, the invariant Hermitian pairing is a
positive definite coefficient pairing tensored with the fixed
rank-six multiplicity pairing.  Restricting the coefficient pairing
to different $K$-lines multiplies the multiplicity pairing by
positive rational numbers.  Thus the new naturally restricted
polarization and the original one differ, under a $K$-isogeny, by
a positive rational scalar $c$.  Its determinant effect is
\[
                         c^6=N_{K/\mathbf Q}(c^3).
 \tag{NID.12}
\]
It leaves the signed determinant class unchanged.  No extension of
$K$ is used for this comparison, and no proper-field norm map is
declared injective.

The period map of the quotient family factors through its actual
Hurwitz component.  Rank nine on the original-family image implies
rank at least nine on that component; its signature-$(3,3)$
unitary period domain has dimension nine.  The natural rational
polarization form, and hence its determinant class, is locally
constant on this component.  It is consequently a valid replacement
source in the same determinant class, not only an isogeny at one
isolated fibre.

We spell out the genus restriction used twice below.  For a nontrivial
irreducible coefficient representation of degree $d\ge2$ and
complex cohomological multiplicity six, over a base of genus $h$,
write $R$ for the genuine branch count and $D$ for the sum of its
moving ranks.  The Euler formula and the finite pointed-base
Hurwitz map give
\[
 6=2d(h-1)+D,
 \qquad R\le D,
 \qquad \operatorname{rank}(\mathrm{period})
       \le\dim\mathcal M_{h,R}.
 \tag{NID.13}
\]
For $h=1$, $D=6$ and the pointed moduli dimension is at most six.
For $h=2$, $D=6-2d\ge0$ forces $d\le3$ and the dimension is at
most $3+D=9-2d\le5$.  For $h\ge3$, $D<0$.
Hence rank nine forces $h=0$, at least twelve genuine branch
points, and $D=2d+6$.  This is the actual-cover argument of
D-BI-01 and D-BI-18; its use does not infer rank from a Hurwitz
space dimension.  Apply it first to $W_K$, and then to $\tau$
using the complex Frobenius identification.  Both $Z=C/H'$ and
$X=C/H$ are therefore rational.

\subsection*{4. The abelian rational map and the index bound}

Assume $\ell>1$.  The map $X\to Z$ is an actual abelian
Galois cover of degree $\ell$ and group $A=H'/H$.
A finite abelian subgroup of $\operatorname{PGL}_2(\mathbf C)$
is cyclic or a Klein four group.  Here is the elementary proof
needed for this instance.  If an element has order greater than
two, any commuting element preserves its two fixed points
individually: interchanging them would invert that element.
The whole group is then a finite subgroup of the multiplicative
group of the coordinate fixing those points, hence cyclic.
Otherwise all nonidentity elements are involutions.  Put one of
them in the form $z\mapsto-z$.  Its centralizer consists of maps
$z\mapsto az$ and $z\mapsto a/z$.  An exponent-two abelian
subgroup contains at most two of each form, so has order at most
four.  The noncyclic possibility is a Klein four group.

In the cyclic case $X\to Z$ has exactly two branch values, both
with full ramification order $\ell$.  In the Klein-four case it
has exactly three branch values with order two.  The latter also
follows directly from genus-zero Riemann--Hurwitz, since every
nonidentity stabilizer has order two: in the displayed coordinate
description, the two fixed points of $z\mapsto-z$ are exchanged
by every element of the other form, and the same holds for each
of the three involutions after conjugation.  Write $s=2$ or $3$
accordingly.

Take the genuine inertia tuple $g_1,\ldots,g_r$ for the reduced
actual $J'$-cover $Y\to Z$.  The preceding section gives
\[
 r\ge12,
 \qquad \sum_{i=1}^r\operatorname{rank}(1-W(g_i))=2e\ell+6.
 \tag{NID.14}
\]
Its images $\bar g_i$ in $A$ are the monodromies of $X\to Z$.
Exactly $s$ are nonidentity.  If $\bar g_i$ has order $d_i>1$,
its action on the $\ell$ constituent blocks is regular, with
$\ell/d_i$ cycles.  The fixed space on each cycle has dimension
at most $e$, even when the product of the linear maps around that
cycle is nontrivial.  Consequently
\[
 \operatorname{rank}(1-W(g_i))\ge e\ell(1-1/d_i),
 \qquad
 \sum_{\bar g_i\ne1}\ell(1-1/d_i)=2\ell-2.
 \tag{NID.15}
\]
The second equality is the actual Riemann--Hurwitz formula for
$X\to Z$.

If $\bar g_i=1$, then $g_i$ is a nonidentity element of $H/N$.
This time every block is nontrivial: $\tau$ is faithful on
$H/N$, and all $\ell$ constituent characters are Galois
conjugate, with the same kernel.  More precisely their moving
ranks are equal, since the multiplicity of the eigenvalue one
is unchanged by a coefficient automorphism.  Thus
$\operatorname{rank}(1-W(g_i))\ge\ell$.
This is the step which would be false for arbitrary different
conjugate kernels before the orbit replacement.

Combining these two kinds of inertia with (NID.14) yields
\[
 2e\ell+6\ge e(2\ell-2)+\ell(r-s),
 \qquad \ell(r-s)\le2e+6.
 \tag{NID.16}
\]
For $s=2$ and $r\ge12$ this is (NID.4).  For $s=3$ and
$\ell=4$ it gives $36\le2e+6$, hence $e\ge15$.
The case $\ell=1$ needs no bound beyond $1\le b(e)$.
This proves (NID.5)--(NID.6).

\subsection*{5. Minimality and the exact boundary}

The $\Gamma$-orbit has $\ell\le k$ elements.  If $\ell<k$,
the replacement coefficient degree $e\ell$ is strictly less
than $ek$.  It belongs to the full comparison class, including
nonprimitive sources, and has the same exact field and natural
determinant.  Minimality therefore forces $\ell=k$ and $H'=G$.
Then (NID.10), together with faithfulness of $\rho_K=W_K$,
gives $N=1$.  No faithfulness of $\tau$ on the original $H$
was used before this conclusion.

The theorem does not make a minimal source complex primitive.
It controls a normal inducing subgroup and permits $L>K$.
It makes no assertion for a nonnormal inducing subgroup:
there the conjugates of $\tau$ need not be a regular Clifford
orbit, and Galois can move a complex block system to a different
block system.  The orbit sum in $K[H]$ then need not give a
transitively permuted direct-sum system for $G$.
The case $e=1$ is not covered by (NID.13)'s rational-base
argument; in particular an unramified genus-four cyclic
intermediate source is not excluded by that argument.
The separate accepted monomial-source results retain their
own scope.  No new rank-nine source, prescribed detector cycle,
or universal original-field descent is asserted here.
